\documentclass[11pt,a4paper]{report}
\usepackage[utf8]{inputenc}
\usepackage[T1]{fontenc}
\usepackage[french]{babel}
\usepackage{lmodern}
\usepackage{geometry}
\usepackage{graphicx}
\usepackage{xcolor}
\usepackage{titlesec}
\usepackage{fancyhdr}
\usepackage{enumitem}
\usepackage{hyperref}
\usepackage{setspace}
\usepackage{array}
\usepackage{tabularx}

\geometry{margin=2.5cm}
\setstretch{1.15}
\pagestyle{fancy}
\fancyhf{}
\fancyfoot[C]{\thepage}
\renewcommand{\headrulewidth}{0pt}

\titleformat{\chapter}[display]
  {\normalfont\Large\bfseries}{\chaptertitlename\ \thechapter}{20pt}{\Huge}

\begin{document}

\begin{titlepage}
    \centering
    \vspace*{2cm}
    {\Large \textbf{Rapport de stage}}\par
    \vspace{1cm}
    {\LARGE \textbf{Développement d’un système de gestion des formations et intégration d’un assistant intelligent avec étude de sécurité}}\par
    \vspace{1.5cm}
    {\large Présenté par :}\par
    \vspace{0.3cm}
    {\Large \textbf{Jmel Ismail}}\par
    \vspace{0.8cm}
    {\large Filière / spécialité}\par
    \vspace{0.2cm}
    {\large \textbf{1ere année cycle ingénieur en informatique}}\par
    \vspace{1.5cm}
    {\large Encadré par :}\par
    \vspace{0.2cm}
    {\large \textbf{Saadaoui Med Ali}}\par
    \vspace{2cm}
    {\large \today}\par
    \vspace{1cm}
    \rule{0.8\textwidth}{0.5pt}
    \vspace{0.5cm}
    {\large ecole nationale des sciences de l’informatique (ENSI) - Tunisie}\par
\end{titlepage}

\tableofcontents
\thispagestyle{empty}
\clearpage

\chapter*{Introduction}
\addcontentsline{toc}{chapter}{Introduction}
Au cours de ce projet, il a été question de mettre en place une solution complète permettant de gérer les inscriptions, les formations, ainsi que l’interaction avec un assistant intelligent destiné à accompagner les utilisateurs. Ce rapport présente les principales étapes réalisées, les changements apportés au système et les résultats obtenus jusqu’à ce niveau de développement.

L’objectif global était de créer une plateforme simple, fonctionnelle et évolutive, capable de répondre aux besoins de gestion administrative tout en intégrant une interface de chatbot pour améliorer l’expérience utilisateur.

\chapter{Présentation du projet}
Le projet consiste en la conception et le développement d’une application web orientée vers la gestion des formations et des participants. Il s’appuie sur une architecture simple composée de pages publiques, d’un espace administrateur et d’un module de chatbot.

Le système permet notamment de:
\begin{itemize}[leftmargin=1.5em]
    \item gérer les cycles de formation ;
    \item enregistrer les participants à une formation ;
    \item administrer les données à travers une interface dédiée ;
    \item proposer une assistance automatique via un chatbot.
\end{itemize}

Les principaux composants du projet sont développés à partir de pages PHP, de fichiers de configuration, d’un widget de chat et d’une intégration d’API pour l’intelligence artificielle.

\chapter{Étapes réalisées}
\section{Analyse et préparation du projet}
La première étape a consisté à analyser les besoins du système ainsi que la structure générale du projet. Il a été nécessaire de clarifier les fonctionnalités attendues, notamment la gestion des inscriptions, la gestion des cycles et la mise à disposition d’une interface de dialogue intelligente.

Cette phase a permis de définir l’organisation du code et de préparer la base du projet en séparant les parties publiques, administratives et techniques.

\section{Mise en place de l’interface de participation}
Une page publique a été conçue pour permettre aux utilisateurs de s’inscrire à une formation. Cette interface rassemble les informations principales demandées, telles que le nom, le prénom, la CIN, l’entreprise, les coordonnées et la formation choisie.

Cette étape a permis d’introduire une première interaction utilisateur avec le système et de poser les bases de la collecte des données.

\section{Développement de la partie administrative}
Une partie administrative a ensuite été mise en place afin de permettre une gestion plus fluide des données. Les formulaires et les interfaces de gestion ont été organisés pour faciliter l’ajout, la modification et la consultation des informations relatives aux cycles, aux formateurs et aux participants.

Le travail réalisé à ce niveau a permis de renforcer la logique de gestion du système et de faire évoluer l’application vers une plateforme plus complète.

\section{Intégration du chatbot}
Une fonctionnalité de chatbot a ensuite été ajoutée afin d’améliorer l’accompagnement des utilisateurs. Cette intégration a été réalisée à travers un widget interactif visible sur l’interface, ainsi qu’un script côté client et un traitement côté serveur.

Le chatbot a été conçu pour répondre à des questions liées aux formations, aux inscriptions et à la navigation du site. La communication avec l’API a été mise en place dans un objectif de rendre l’assistant plus utile et plus réactif.

\section{Configuration et tests}
La dernière étape a porté sur la configuration de l’accessibilité de l’assistant intelligent. Dans cette version, l’approche a évolué vers un modèle plus sécurisé en privilégiant une solution locale plutôt que l’utilisation d’une clé API externe. Cette évolution a permis de simplifier l’intégration tout en renforçant la confidentialité des échanges.

Des vérifications ont également été effectuées afin de garantir le bon fonctionnement des différents modules, notamment l’affichage du widget, l’envoi des messages et la réponse du système.

\chapter{Étude de sécurité des données et migration vers Ollama}
L’un des aspects essentiels du projet concerne la sécurité des données des utilisateurs ainsi que la protection des informations transmises à l’agent intelligent. Les données traitées incluent des informations personnelles sensibles, telles que le nom, la CIN, l’entreprise, les coordonnées téléphoniques et l’adresse électronique. Il est donc indispensable de garantir leur confidentialité, leur intégrité et leur protection contre toute fuite ou utilisation non autorisée.

\section{Problématique de sécurité}
L’utilisation d’une clé API externe présente plusieurs risques potentiels. D’une part, elle nécessite l’envoi des requêtes vers un service tiers, ce qui peut exposer certaines données à des environnements extérieurs non maîtrisés. D’autre part, la gestion d’une clé API implique des précautions particulières pour éviter toute fuite, toute mauvaise configuration ou toute réutilisation non intentionnelle. Enfin, cette approche dépend fortement de la disponibilité du service externe ainsi que de la qualité de la connexion réseau.

\section{Choix stratégique : utiliser Ollama}
Face à ces contraintes, il a été décidé de ne plus s’appuyer sur une clé API et d’opter pour une solution basée sur Ollama. Ollama permet d’exécuter des modèles de langage localement, ce qui réduit considérablement les risques liés à l’exposition des données à des serveurs distants. Cette approche offre un meilleur contrôle sur les informations traitées et renforce la confidentialité des échanges entre l’utilisateur et l’agent IA.

Le passage à Ollama présente plusieurs avantages :
\begin{itemize}[leftmargin=1.5em]
    \item protection accrue des données sensibles ;
    \item réduction du risque de fuite liée à la clé API ;
    \item meilleure maîtrise du traitement local des informations ;
    \item indépendance vis-à-vis des services cloud externes ;
    \item conformité plus adaptée à un contexte où la sécurité des données est prioritaire.
\end{itemize}

\section{Mesures de sécurité recommandées}
Même avec une solution locale comme Ollama, certaines mesures de sécurité restent indispensables. Il est recommandé de :
\begin{itemize}[leftmargin=1.5em]
    \item limiter l’accès au serveur local aux utilisateurs autorisés ;
    \item sécuriser les environnements d’hébergement et les droits d’accès ;
    \item éviter le stockage inutile des conversations sensibles ;
    \item mettre à jour régulièrement les composants du système ;
    \item surveiller les logs et les accès afin de détecter toute activité suspecte.
\end{itemize}

\section{Conclusion partielle}
Le choix d’Ollama constitue une solution plus respectueuse de la confidentialité et plus adaptée aux exigences de sécurité du projet. Il permet de concilier simplicité d’utilisation, protection des données et meilleure autonomie technologique, tout en répondant aux besoins fonctionnels du système.

\section{Configuration du script PHP de communication avec Ollama}
Une étape importante de l’évolution du système a consisté à mettre en place un script PHP chargé d’interagir avec l’agent IA local. Ce script joue le rôle d’intermédiaire entre l’interface utilisateur et le moteur d’intelligence artificielle exécuté localement via Ollama.

Son fonctionnement repose sur plusieurs étapes essentielles. Tout d’abord, il reçoit le message envoyé par l’utilisateur à travers l’interface JavaScript. Ensuite, il prépare une requête vers le serveur local d’Ollama, disponible sur le port 11434, sans nécessiter de clé API externe. L’IA est ensuite sollicitée avec un message système défini au préalable, afin de lui donner un rôle clair et adapté au contexte du centre de formation. Enfin, la réponse générée par le modèle est récupérée et renvoyée vers l’interface utilisateur sous forme de réponse exploitable.

Cette solution présente un intérêt majeur, car elle permet de conserver un fonctionnement local, simple et sécurisé, tout en rendant possible l’intégration d’un assistant conversationnel dans le projet. Le script PHP assure ainsi la liaison entre la logique métier du site et le moteur IA, dans un cadre plus maîtrisé et plus conforme aux exigences de confidentialité.

\chapter{Changements majeurs apportés}
Au cours de l’évolution du projet, plusieurs changements importants ont été réalisés :

\begin{itemize}[leftmargin=1.5em]
    \item création d’une page d’accueil publique dédiée à l’inscription ;
    \item ajout d’un espace administratif avec des formulaires de gestion ;
    \item mise en place d’un système de navigation entre les différentes pages administratives ;
    \item intégration d’un widget de chatbot dans l’interface ;
    \item ajout d’une configuration de clé API pour l’activation du chatbot ;
    \item développement d’un traitement backend pour envoyer les messages à l’API et récupérer les réponses.
\end{itemize}

Ces changements ont permis de transformer le projet d’une simple structure statique vers une application plus dynamique et plus interactive.

\chapter{Résultats obtenus}
À ce stade, le projet présente une structure fonctionnelle et cohérente. L’utilisateur peut accéder à une interface d’inscription, l’administrateur peut gérer les informations principales, et un assistant intelligent est désormais intégré pour répondre à certaines requêtes.

Le système offre ainsi un cadre plus pratique pour la gestion des formations, tout en améliorant l’expérience utilisateur grâce à l’automatisation et à la présence d’un assistant conversationnel.

\chapter{Conclusion et perspectives}
Ce projet a permis de mettre en place une base solide pour un système de gestion des formations enrichi par une intelligence artificielle. Les étapes réalisées ont permis de couvrir les besoins fondamentaux de gestion, d’interaction et de personnalisation.

Les perspectives d’évolution sont nombreuses. Il serait intéressant, à court terme, d’améliorer la qualité des réponses du chatbot, d’ajouter davantage de fonctionnalités de gestion, et d’optimiser l’interface pour la rendre plus fluide et plus moderne.

\chapter*{Conclusion générale}
\addcontentsline{toc}{chapter}{Conclusion générale}
En somme, le travail accompli jusqu’à présent constitue une première version fonctionnelle du système. Il offre une bonne base technique et fonctionnelle, et il ouvre la voie à des améliorations plus avancées et plus professionnelles.

\end{document}
